\documentclass[../../main.tex]{subfiles}
\begin{document}

{\bf [解]}
条件(1)より直線$L$が$A(0,1)$を通るから，この$L$上の点はベクトル表記で
\begin{align}
 \begin{pmatrix}
  x \\
  y
 \end{pmatrix}  
 = 
 \begin{pmatrix}
  0 \\
  1
 \end{pmatrix}
 + t 
 \begin{pmatrix}
  \cos \theta \\
  \sin \theta
 \end{pmatrix} \label{eq:1}
\end{align}
と書ける．ただし$t$は実数，$\theta$は
\begin{align}
 -\dfrac{\pi}{2} < \theta \le \dfrac{\pi}{2} \label{eq:2}
\end{align}
とする．

さて，$L$上の点$Q$が$t=t_1$に対応するとすると，条件(2)より$f$による写像$P=f(Q)$：
\begin{align}
 f(Q) =
  \begin{pmatrix}
   at_1\cos\theta + (a-2)\left(t_1\sin\theta + 1\right) \\
   (a-2)t_1\cos\theta + a\left(t_1\sin\theta + 1\right)
  \end{pmatrix}
\end{align}
も$L$上にある．従ってベクトル$\overrightarrow{AP}$が$L$の方向ベクトル$(\cos\theta,\sin\theta)$と並行であれば良い．
実際計算すると
\begin{align}
 \overrightarrow{AP} = 
  \begin{pmatrix}
   \left(a\cos\theta + (a-2)\sin\theta\right) t_1 + (a-2) \\
   \left(a\sin\theta + (a-2)\cos\theta\right) t_1 + (a-1) 
  \end{pmatrix} \label{eq:3}
\end{align}
だから，ベクトルの並行条件より
\begin{align}
 & \left[\left(a\sin\theta + (a-2)\cos\theta\right) t_1 + (a-1)\right]\cos\theta - \left[\left(a\cos\theta + (a-2)\sin\theta\right) t_1 + (a-2)\right]\sin\theta = 0 \\
 \left[a\cos\theta\sin\theta + (a-2)\sin^2\theta -a\cos\theta\sin\theta - (a-2)\cos^2\theta\right]t_1 + (a-2)\sin\theta - (a-1)\cos\theta = 0 \\
 (a-2)\left(\sin^2\theta -\cos^2\theta\right)t_1 + \left[(a-1)\sin\theta-(a-2)\cos\theta\right] = 0 \label{eq:4}
\end{align}
である．これが$t_1$によらず成立する，すなわち$t_1$についての恒等式だから係数が全て$0$であればよく，
\begin{align}
 \begin{cases}
  (a-2)\left(\sin^2\theta -\cos^2\theta\right) = 0 \\
  (a-1)\sin\theta-(a-2)\cos\theta = 0
 \end{cases} \label{eq:5}
\end{align}
を得る．\cref{eq:5}の第一式より$a-2=0$または$\sin^2\theta=\cos^2\theta$が必要である．

\subparagraph{$a=2$の時}
この時\cref{eq:5}の第二式は$\sin\theta=0$となる．これを満たすのは\cref{eq:2}の範囲では$\theta=0$の時であり十分である．

\subparagraph{$\sin^2\theta=\cos^2\theta$の時}

この時は$\sin\theta=\pm \cos\theta$であるから\cref{eq:5}の第二式に代入すると
\begin{align}
 \pm(a-1)\cos\theta = (a-2)\cos\theta \\
\therefore
 \begin{cases}
  \cos\theta = 0 & (\text{複合正}) \\
  (2a-3)\cos\theta = 0 & (\text{複合負})
 \end{cases}
\end{align}
である．$\cos\theta=0$の時は$\sin\theta=\pm\cos\theta=0$となるが，$\sin\theta$と$\cos\theta$を同時に$0$にする$\theta$は存在しないため不適で，従って$a=\dfrac{3}{2}$の時のみ条件を満たして十分である．
\casesend

以上から条件をみたす$a$は
\begin{align}
a=2,\ \dfrac{3}{2}
\end{align}
である．


\newpage

\end{document}
